\subsection{Shared research guidance and evaluation calibration}
\label{app:hitl-protocol}
The controlled study uses Gemini 3.5 Flash through Vertex
\texttt{generateContent}, with a local program-development controller.
It is distinct from the managed AlphaEvolve discovery campaigns.
The two policies receive identical instructions to maintain and mutate
$L=S^\perp$ throughout search and convert candidate bases to $S$ using
binary nullspace and symplectic-swap helpers. The initial code is the
same $S$-side annealing callback in both policies. There is no protected
initializer. Programs may change representations, objectives, mutation
operators and restart schedules; the working $L$-side calibration code,
archived codes and family formula are not provided to either policy.

Before model generation, four calibration seeds compare the initial
callback with an interface-adapted, researcher-assisted $L$-side control
at up to $10^5$ evaluator calls. The controls depend on the call limit
only through the remaining allowance, so their first-hit records give
success curves at 1,000, 3,000, 10,000, 30,000 and 100,000 calls.
At $n=9$, the rule selects the smallest grid budget with initial
success between 0.1 and 0.5, control success at least 0.5 and a difference
of at least 0.25. A second cell may have zero initial success if the
control succeeds in at least half the executions. The selected budget
is 10,000: initial/control success is 1/4 versus 2/4 at $n=9$, and
0/4 versus 2/4 at $n=11$. Calibration is a test of whether the evaluator
can reward a useful edit, not a measurement of evolutionary contribution.
Independently verified codes establish feasibility at all training and transfer
parameters before generation.

\subsection{Generation, interruption and selection}\label{app:hitl-selection}
Four blocks each plan 12 proposals per policy. Independent proposals
receive the fixed initial source and its fixed training feedback.
Iterative evolution receives the best-so-far source and the latest
attempt, including rejected attempts, with their training feedback.
Training uses two fixed execution seeds at $n=9,11$. Each proposal is
ranked lexicographically by mean normalized success-curve area, success
fraction and clipped distance ratio; earlier proposals break remaining
ties. For first-hit call $h$ under allowance $B$, the primary utility is
\[
 u(h;B)=\begin{cases}1-h/B,&\text{if the target is solved by call }B,\\
 0,&\text{otherwise.}\end{cases}
\]
This equals the normalized integral of the success indicator over the
available evaluator-call budget. The initial program remains eligible
for selection. Proposal counts match; token and execution costs need not.

The planned 12-proposal runs received 58 model responses before
unavailable responses interrupted generation. Before any validation or
test observations, a completion-only rule retained the longest fully
received prefix shared by every policy and block: seven proposals each,
56 total. The two additional responses and all request records are
preserved separately. Earlier exploratory studies informed calibration
and guidance; their measurements are not pooled into this comparison.

For the included 28 requests per policy, recorded input/output/reasoning
token totals are 50,344/22,891/22,780 for independent proposals and
86,501/26,145/32,128 for iterative evolution. Mean request latency is
8.77\,s and 11.93\,s, respectively. The additional source and feedback
in iterative prompts therefore incur greater model usage despite equal
proposal counts. Receipt-level accounting is preserved in the artifact; service latency
is separate from search time.

After generation ends, four new validation seeds at the training
parameters select among each policy/block's three training-best distinct
programs, using the same metric order and then training rank. The eight
selected programs are distinct. Test uses eight further seeds at each
training parameter and 10,000 calls per execution. The separate length
transfer uses $n=13,15$, four seeds and 30,000 calls. Transfer measurements
never enter selection. All seed lists, requests, source hashes, parent
relationships, candidate outcomes and selections are retained in
\texttt{experiments/hitl\_ablation/}.

% Start the results with their table rather than leaving a heading at the page foot.
\subsection{Complete policy outcomes}\label{app:hitl-outcomes}
Table~\ref{tab:hitl-blocks} reports the original test cohort and the
32-seed repeat of the same validation-selected programs. Displayed values
are rounded; differences use unrounded values. The new seeds
were fixed before execution and did not change any program or selection.
The original paired mean curve-area difference is $+0.1188$ (two positive
blocks), versus $+0.0044$ in the repeat (one positive block). Descriptive
two-sided sign-flip values over the four block differences are 0.50 and
1.00. These calculations use four generation blocks; execution seeds
within one program quantify its run-to-run variability. Excluding the
block affected by later unavailable responses gives an original mean
difference of $+0.0819$. All included proposals precede those responses.

\begin{table}[!htbp]
\centering\small
\begin{tabular}{crrrrrr}
\toprule
& \multicolumn{2}{c}{Original /16} & \multicolumn{2}{c}{Repeat /64}
& \multicolumn{2}{c}{Curve-area difference}\\
\cmidrule(lr){2-3}\cmidrule(lr){4-5}\cmidrule(lr){6-7}
Block & Independent & Iterative & Independent & Iterative & Original & Repeat\\
\midrule
1 & 4 & 3 & 38 & 7 & $-0.114$ & $-0.201$ \\
2 & 4 & 10 & 27 & 26 & $+0.229$ & $-0.014$ \\
3 & 5 & 16 & 34 & 49 & $+0.395$ & $+0.287$ \\
4 & 8 & 9 & 36 & 33 & $-0.036$ & $-0.055$ \\
\midrule
Total/mean & 21/64 & 38/64 & 135/256 & 115/256 & $+0.119$ & $+0.004$\\
\bottomrule
\end{tabular}
\caption{Test outcomes by generation block.}\label{tab:hitl-blocks}
\end{table}

Iterative selections are proposals 2, 1, 3 and 1 in blocks 1--4,
respectively (proposal numbers start at one). Thus two selections are
first proposals and two incorporate feedback from earlier generated
programs. The classical $L$-side control uses the shared proposed
representation with the preserved annealing objective, moves and
restart policy; its source is unchanged from calibration.

The repeat covers 32 new seeds at each of $n=9,11,13,15$, at 10,000
calls for the first two lengths and 30,000 for the latter two. It
includes all eight selected programs, both fixed controls, and the
first two programs in the block-3 iterative lineage. Together with
24 executions of the classical control on the original test and
longer-code seeds, this gives 1,560 executions and 545 independently
verified successful code records, with no process errors. All records,
program hashes and the frozen protocol are in
\texttt{experiments/hitl\_followup/}.

At $n=9,11$, the repeat gives 96/128 and 39/128 successes for independent
proposals, versus 81/128 and 34/128 for iterative evolution. The classical
control gives 26/32 and 10/32, and the initial program 1/32 and 0/32.
Figure~\ref{fig:hitl-followup} shows the longer-code curves; the artifact
also retains the complete four-length plot.

\begin{figure}[!htbp]
\centering
\resizebox{\linewidth}{!}{\definecolor{hfinitial}{HTML}{A37545}
\definecolor{hfind}{HTML}{78818E}
\definecolor{hfevo}{HTML}{00858B}
\definecolor{hfcontrol}{HTML}{B55D32}
\begin{tikzpicture}[x=1cm,y=1cm,font=\scriptsize]
\path[use as bounding box] (-1,-1.70) rectangle (13.3,3.7);
\begin{scope}[xshift=0cm,yshift=0cm]
\node[font=\scriptsize\bfseries] at (2.825,3.2) {(a) $n=13,\ d=12$};
\draw[black!12,line width=.35pt] (0,0.0)--(5.65,0.0);
\node[anchor=east] at (-.12,0.0) {0};
\draw[black!12,line width=.35pt] (0,0.7)--(5.65,0.7);
\node[anchor=east] at (-.12,0.7) {25};
\draw[black!12,line width=.35pt] (0,1.4)--(5.65,1.4);
\node[anchor=east] at (-.12,1.4) {50};
\draw[black!12,line width=.35pt] (0,2.1)--(5.65,2.1);
\node[anchor=east] at (-.12,2.1) {75};
\draw[black!12,line width=.35pt] (0,2.8)--(5.65,2.8);
\node[anchor=east] at (-.12,2.8) {100};
\draw[black!40] (0,2.8)--(0,0)--(5.65,0);
\node[rotate=90] at (-.72,1.4) {Success rate (\%)};
\node at (2.825,-.60) {Exact evaluations};
\node[anchor=north] at (0.0,-.10) {0};
\node[anchor=north] at (1.8833333333333333,-.10) {10k};
\node[anchor=north] at (3.7666666666666666,-.10) {20k};
\node[anchor=north] at (5.65,-.10) {30k};
\draw[hfind,line width=1.0pt] (0.00000,0.00000) -- (1.07237,0.00000) -- (1.07237,0.02187) -- (1.30138,0.02187) -- (1.30138,0.04375) -- (1.46787,0.04375) -- (1.46787,0.06562) -- (1.76525,0.06562) -- (1.76525,0.08750) -- (3.14385,0.08750) -- (3.14385,0.10938) -- (3.45422,0.10938) -- (3.45422,0.13125) -- (3.87496,0.13125) -- (3.87496,0.15312) -- (3.91488,0.15312) -- (3.91488,0.17500) -- (4.14107,0.17500) -- (4.14107,0.19687) -- (4.55823,0.19687) -- (4.55823,0.21875) -- (4.55899,0.21875) -- (4.55899,0.24062) -- (4.68121,0.24062) -- (4.68121,0.26250) -- (4.70739,0.26250) -- (4.70739,0.28437) -- (4.73602,0.28437) -- (4.73602,0.30625) -- (5.46751,0.30625) -- (5.46751,0.32812) -- (5.50442,0.32812) -- (5.50442,0.35000) -- (5.61007,0.35000) -- (5.61007,0.37187) -- (5.65000,0.37187);
\draw[hfevo,line width=1.0pt] (0.00000,0.00000) -- (0.29964,0.00000) -- (0.29964,0.02187) -- (0.80776,0.02187) -- (0.80776,0.04375) -- (0.82377,0.04375) -- (0.82377,0.06562) -- (0.82528,0.06562) -- (0.82528,0.08750) -- (0.98404,0.08750) -- (0.98404,0.10938) -- (1.04808,0.10938) -- (1.04808,0.13125) -- (1.14883,0.13125) -- (1.14883,0.15312) -- (1.51194,0.15312) -- (1.51194,0.17500) -- (1.57334,0.17500) -- (1.57334,0.19687) -- (1.59236,0.19687) -- (1.59236,0.21875) -- (1.62362,0.21875) -- (1.62362,0.24062) -- (1.72702,0.24062) -- (1.72702,0.26250) -- (1.87900,0.26250) -- (1.87900,0.28437) -- (1.90763,0.28437) -- (1.90763,0.30625) -- (1.94153,0.30625) -- (1.94153,0.32812) -- (2.08598,0.32812) -- (2.08598,0.35000) -- (2.12120,0.35000) -- (2.12120,0.37187) -- (2.16979,0.37187) -- (2.16979,0.39375) -- (2.17299,0.39375) -- (2.17299,0.41562) -- (2.29691,0.41562) -- (2.29691,0.43750) -- (2.42310,0.43750) -- (2.42310,0.45937) -- (2.46867,0.45937) -- (2.46867,0.48125) -- (2.49353,0.48125) -- (2.49353,0.50312) -- (2.75833,0.50312) -- (2.75833,0.52500) -- (2.91069,0.52500) -- (2.91069,0.54688) -- (3.25986,0.54688) -- (3.25986,0.56875) -- (3.42955,0.56875) -- (3.42955,0.59062) -- (3.44951,0.59062) -- (3.44951,0.61250) -- (3.55121,0.61250) -- (3.55121,0.63437) -- (3.91074,0.63437) -- (3.91074,0.65625) -- (4.00623,0.65625) -- (4.00623,0.67812) -- (4.17780,0.67812) -- (4.17780,0.70000) -- (4.40399,0.70000) -- (4.40399,0.72187) -- (4.54147,0.72187) -- (4.54147,0.74375) -- (4.66728,0.74375) -- (4.66728,0.76562) -- (5.04790,0.76562) -- (5.04790,0.78750) -- (5.28840,0.78750) -- (5.28840,0.80937) -- (5.45507,0.80937) -- (5.45507,0.83125) -- (5.49594,0.83125) -- (5.49594,0.85312) -- (5.65000,0.85312);
\draw[hfcontrol,line width=1.0pt,dash pattern=on 4pt off 1.5pt] (0.00000,0.00000) -- (1.46787,0.00000) -- (1.46787,0.08750) -- (4.70739,0.08750) -- (4.70739,0.17500) -- (5.61007,0.17500) -- (5.61007,0.26250) -- (5.65000,0.26250);
\draw[hfinitial,line width=1.0pt,dash pattern=on 3pt off 2pt] (0.00000,0.00000) -- (5.65000,0.00000);
\end{scope}
\begin{scope}[xshift=7.2cm,yshift=0cm]
\node[font=\scriptsize\bfseries] at (2.825,3.2) {(b) $n=15,\ d=14$};
\draw[black!12,line width=.35pt] (0,0.0)--(5.65,0.0);
\node[anchor=east] at (-.12,0.0) {0};
\draw[black!12,line width=.35pt] (0,0.7)--(5.65,0.7);
\node[anchor=east] at (-.12,0.7) {25};
\draw[black!12,line width=.35pt] (0,1.4)--(5.65,1.4);
\node[anchor=east] at (-.12,1.4) {50};
\draw[black!12,line width=.35pt] (0,2.1)--(5.65,2.1);
\node[anchor=east] at (-.12,2.1) {75};
\draw[black!12,line width=.35pt] (0,2.8)--(5.65,2.8);
\node[anchor=east] at (-.12,2.8) {100};
\draw[black!40] (0,2.8)--(0,0)--(5.65,0);
\node[rotate=90] at (-.72,1.4) {Success rate (\%)};
\node at (2.825,-.60) {Exact evaluations};
\node[anchor=north] at (0.0,-.10) {0};
\node[anchor=north] at (1.8833333333333333,-.10) {10k};
\node[anchor=north] at (3.7666666666666666,-.10) {20k};
\node[anchor=north] at (5.65,-.10) {30k};
\draw[hfind,line width=1.0pt] (0.00000,0.00000) -- (3.09639,0.00000) -- (3.09639,0.02187) -- (5.65000,0.02187);
\draw[hfevo,line width=1.0pt] (0.00000,0.00000) -- (0.57065,0.00000) -- (0.57065,0.02187) -- (0.63337,0.02187) -- (0.63337,0.04375) -- (0.66953,0.04375) -- (0.66953,0.06562) -- (1.17897,0.06562) -- (1.17897,0.08750) -- (1.39103,0.08750) -- (1.39103,0.10938) -- (1.43359,0.10938) -- (1.43359,0.13125) -- (1.44471,0.13125) -- (1.44471,0.15312) -- (1.66204,0.15312) -- (1.66204,0.17500) -- (1.68408,0.17500) -- (1.68408,0.19687) -- (2.08447,0.19687) -- (2.08447,0.21875) -- (2.66850,0.21875) -- (2.66850,0.24062) -- (2.67716,0.24062) -- (2.67716,0.26250) -- (3.29903,0.26250) -- (3.29903,0.28437) -- (4.24880,0.28437) -- (4.24880,0.30625) -- (4.50550,0.30625) -- (4.50550,0.32812) -- (4.68366,0.32812) -- (4.68366,0.35000) -- (4.77067,0.35000) -- (4.77067,0.37187) -- (5.29876,0.37187) -- (5.29876,0.39375) -- (5.34471,0.39375) -- (5.34471,0.41562) -- (5.40328,0.41562) -- (5.40328,0.43750) -- (5.65000,0.43750);
\draw[hfcontrol,line width=1.0pt,dash pattern=on 4pt off 1.5pt] (0.00000,0.00000) -- (5.65000,0.00000);
\draw[hfinitial,line width=1.0pt,dash pattern=on 3pt off 2pt] (0.00000,0.00000) -- (5.65000,0.00000);
\end{scope}
\draw[hfinitial,line width=1pt,dash pattern=on 3pt off 2pt] (1,-1.1)--(1.55,-1.1);
\node[anchor=west] at (1.65,-1.1) {Initial program};
\draw[hfcontrol,line width=1pt,dash pattern=on 3pt off 2pt] (7.1,-1.1)--(7.6499999999999995,-1.1);
\node[anchor=west] at (7.75,-1.1) {Classical L-side control};
\draw[hfind,line width=1pt] (1,-1.53)--(1.55,-1.53);
\node[anchor=west] at (1.65,-1.53) {Independent proposals};
\draw[hfevo,line width=1pt] (7.1,-1.53)--(7.6499999999999995,-1.53);
\node[anchor=west] at (7.75,-1.53) {Iterative evolution};
\end{tikzpicture}
}
\caption{Longer-code replay on fresh execution seeds.}
\label{fig:hitl-followup}
\end{figure}

\FloatBarrier

Helper computations remain outside the evaluator-call count but inside
the common resource limits. Each target/seed runs in a separate sandbox
with 120/125\,s soft/hard CPU limits, 135\,s wall time and 1\,GiB resident
memory. Every explicit or returned-candidate verification consumes the
same call quota. Execution times and model request latencies are
recorded separately in the artifact.

The artifact audit reconstructs the original proposal inputs and parent
relationships, checks validation selections and independently verifies
saved code records. The complete inputs and replay commands are available
in the \artifactlink{}.

\subsection{Prospective generation replication}\label{app:replication}
The new cohort fixes eight paired blocks and seven proposal slots per
policy before generation. Model, shared task, starting program,
representation helpers, ranking and validation selection match the
original study. New model seeds and disjoint training, validation,
test and longer-code execution seeds separate the cohorts. Each block
uses two training seeds, four validation seeds and 32 held-out seeds
per length. The test budget is 10,000 calls at $n=9,11$; longer-code
performance at 30,000 calls is secondary.

Of 112 allocated slots, 111 return programs: 56 independent and 55
iterative. One response in block 8 is unavailable. Before validation
or testing, the recorded rule retained all eight blocks, using received
programs without replacing the missing proposal. Subsequent iterative
requests use the most recent received source and feedback. The seven
fully received paired blocks provide a sensitivity analysis. No
program is reselected using test or longer-code outcomes.

Table~\ref{tab:replication-blocks} reports all blocks; each longer-length
entry has 32 executions per policy and block. The mean primary
curve-area difference is $+0.0987$; five differences are positive and
the descriptive two-sided sign-flip value is $p=0.3281$. On the seven
complete blocks, the mean difference is $+0.1504$. These are eight
paired generation observations, not 512 independent generation trials.

\begin{table}[!htbp]
\centering\small
\setlength{\tabcolsep}{5pt}
\begin{tabular}{crrrrr}
\toprule
Block & Independent & Iterative & $\Delta$ area & $n=13$ & $n=15$\\
 & \multicolumn{2}{c}{Test successes /64} & & \multicolumn{2}{c}{Independent / iterative}\\
\midrule
1 & 25 & 19 & -0.063 & 3 / 5 & 0 / 2\\
2 & 25 & 30 & +0.007 & 3 / 10 & 0 / 5\\
3 & 24 & 35 & -0.001 & 2 / 15 & 1 / 4\\
4 & 39 & 45 & +0.021 & 17 / 16 & 5 / 7\\
5 & 26 & 64 & +0.578 & 6 / 32 & 2 / 32\\
6 & 34 & 40 & +0.095 & 15 / 8 & 0 / 3\\
7 & 25 & 62 & +0.417 & 3 / 32 & 0 / 28\\
8 & 25 & 2 & -0.264 & 1 / 0 & 0 / 0\\
\midrule
Total/mean & 223/512 & 297/512 & $+0.099$ & 50 / 118 & 8 / 81\\
\bottomrule
\end{tabular}
\caption{New generation blocks.}\label{tab:replication-blocks}
\end{table}


All 16 selected sources are distinct. Source inspection finds sustained
search in $L$, with exact distance feedback obtained through the common
charged evaluator. Rank checks and basis conversions do not compute
distance privately. One program performs uncharged basis-row additions
that leave the subspace unchanged; other proposals evaluate their
candidate subspaces. Reads of the incumbent before search return no
matrix because controller initialization is disabled. The artifact audit
reconstructs all 111 prompts and parent choices, reproduces validation
selection, checks evaluation quotas and independently verifies saved codes.

The received independent requests use 100,632 input, 41,801 output and
58,067 reported reasoning tokens; iterative requests use 165,819,
49,296 and 64,426, respectively. Equal proposal allocations therefore
do not imply equal token costs. Full transport records and the
missing-response rule are preserved in the artifact.

For Figure~\ref{fig:ablation-comparison}, the original initial program
and classical $L$-side control are also run on these execution seeds.
This fixed-control replay was added after the new policy results were
available; neither control was modified, tuned or selected using them.
Each control has 32 runs per length, and is separate from the eight
paired policy observations. These runs contextualize the success curves;
the prespecified policy endpoint and its analysis remain unchanged.


\subsection{Preserved program and generalization to longer codes}\label{app:hitl-program}
In the original four-block cohort, the program selected in block 3 is
\texttt{b02\_evolution\_g02}, whose preserved parent is
\texttt{b02\_evolution\_g01}; generation indices start at zero.
The parent already searches $L$ and penalizes distance shortfall.
The child increases the ebit-mismatch weight from 5,000 to 10,000 and
the offending weight from 1 to 100. It also mixes row replacement,
Pauli perturbations and basis-row mixing, scales mutation size with
length, and restores a saved local best with a capped temperature
adjustment after stagnation. Tinted expressions in Algorithm~\ref{alg:ablation-search} mark these edits
relative to this immediate parent; the inherited $L$-side representation
and distance-shortfall term are not tinted. These changes are observed
in the preserved source, rather than inferred from performance alone. The repeat below measures the combined changes in this final step.

% Source: b02_evolution_g02; shaded edits are relative to b02_evolution_g01.
\begin{algorithm}[!t]
\caption{Evolved search in the symplectic dual.}
\label{alg:ablation-search}
\small
\begin{algorithmic}[1]
\Require target $(n,k_t,c_t,d_t)$, exact-evaluation budget $N$
\Statex \(\Phi(R)=\algdelta{\(10^4\)}|R.c-c_t|+\algdelta{\(100\)}R.\mathrm{offending}+10^3\Delta_d(R)\)
  \algside{Increase penalties for entanglement mismatch and low-weight violations.}
\State $r\gets n+k_t-c_t$
\While{budget remains}
  \State $L\gets$ a random independent list of $r$ Pauli vectors
  \State $C\gets\Phi(\textsc{Eval}(L^\perp))$; \algdelta{$(L_*,C_*)\gets(L,C)$}
    \algside{Retain the best state within each restart.}
  \State $\tau\gets\algdelta{\(50\)}$; $f\gets0$; \algdelta{$g\gets0$}
  \While{$f<\algdelta{\(600\)}$ \textbf{and} budget remains}
    \State $\tau\gets\max(\algdelta{\(0.01\)},0.995\tau)$
    \State $L'\gets\algdelta{\(\textsc{MixedMove}(L)\)}$
      \algside{Add row mixing and allow larger Pauli changes at longer lengths.}
    \State \textbf{if} $\operatorname{rank}\,L'\ne r$ \textbf{then} $f\gets f+1$; \textbf{continue}
    \State $C'\gets\Phi(\textsc{Eval}((L')^\perp))$; $\delta\gets C'-C$
    \If{$\delta\le0$ \textbf{or} $U(0,1)<\exp(-\delta/\tau)$}
      \State $(L,C)\gets(L',C')$
      \If{\algdelta{$C<C_*$}}
        \State \algdelta{$(L_*,C_*)\gets(L,C)$; $g\gets0$}
      \Else
        \State \algdelta{$g\gets g+1$}
      \EndIf
      \State \textbf{if} $\delta<0$ \textbf{then} $f\gets0$ \textbf{else} $f\gets f+1$
    \Else
      \State $f\gets f+1$; \algdelta{$g\gets g+1$}
    \EndIf
    \If{\algdelta{$g>80$}}
      \algside{On stagnation, return to the best state and adjust the temperature.}
      \State \algdelta{$(L,C)\gets(L_*,C_*)$; $g\gets0$}
      \State \algdelta{$\tau\gets\min(30,1.5\tau)$}
    \EndIf
  \EndWhile
\EndWhile
\State \Return the evaluator's best verified record, if any
\end{algorithmic}
\vspace{2pt}
{\footnotesize\algdelta{Shading}: changes from the immediate parent; exact definitions in Appendix~\ref{app:hitl-program}.}
\end{algorithm}


The program cost in Algorithm~\ref{alg:ablation-search} is
\[
 \Phi(R)=10^4|R.c-c_t|+100\,R.\mathrm{offending}+10^3\Delta_d(R),
\]
where $\Delta_d(R)=\max(0,d_t-R.d)$, or $d_t$ if distance is unavailable.
If the offending count or entanglement cost is missing, the score is
$10^9$. These definitions match the preserved source.

In Algorithm~\ref{alg:ablation-search}, a random Pauli perturbation has weight
uniformly sampled from $1,\ldots,\max(2,\lfloor n/2\rfloor)-1$.
\textsc{MixedMove} chooses a row uniformly: with probability 0.3 it
replaces the row by such a perturbation, with probability 0.3 it XORs
the perturbation into the row, and otherwise it samples a second row.
If that row is different, it XORs the two rows and, with probability
0.5, adds a Pauli perturbation. A rank-deficient proposal increments
$f$ and skips evaluation. Otherwise, $f$ resets only on an accepted
improvement over the current cost, and increments on all other moves.
The counter $g$ resets on a new restart-local best and increments on
other evaluated moves. At $g>80$, the program restores that best,
sets $g=0$ and applies $\tau\leftarrow\min(30,1.5\tau)$. This adjustment
can lower a temperature still above 30; it is not always a temperature
increase. The source and its budget behavior remain unchanged.

% Keep the numerical outcomes together below the algorithm.
\newpage
The original child achieves 16/16 primary-test successes and 3/4 at
each longer length. Its successful longer-code executions use
3,260--9,755 calls at $n=13$ and 7,125--16,756 at $n=15$.
Table~\ref{tab:hitl-lineage} reports every preserved lineage stage on
the 32 new seeds; the original validation selection is unchanged.
Between the first proposal and parent, the only functional search change
is the distance-shortfall penalty. Between parent and selected child,
several policies change together. This lineage was chosen for diagnosis
after its original performance, so the table describes that program's
development separately from the four-block policy comparison.

\begin{table}[!htbp]
\centering\small
\begin{tabular}{lrrrr}
\toprule
Program & $n=9$ & $n=11$ & $n=13$ & $n=15$\\
\midrule
First proposal & 27/32 & 19/32 & 13/32 & 3/32\\
Parent: add distance penalty & 28/32 & 23/32 & 29/32 & 28/32\\
Validation-selected child & 30/32 & 19/32 & 30/32 & 20/32\\
\bottomrule
\end{tabular}
\caption{Preserved lineage on fresh execution seeds.}\label{tab:hitl-lineage}
\end{table}

All eight selected programs implement sustained $L$-side repair. Each
passes candidate matrices to the common evaluator throughout search;
none substitutes a private distance test that filters solutions before
their first charged evaluation.
